\documentclass[10pt,letterpaper]{article}
\usepackage{cornell-notes}

\title{Chapter 16: Permutation Cycles}
\author{}
\date{}

\setDocTitle{Chapter 16: Permutation Cycles}
\setDocSubtitle{Combinatorial Algorithms}
\setDocVersion{v1.0}
\setDocStatus{Study Companion}
\setDocOwner{Combinatorial Algorithms Cornell Notes}
\setDocAuthor{}
\setDocDate{\today}
\setDocSummary{Cornell-format study and review notes for Chapter 16: Permutation Cycles.}

\setCornellCollection{Combinatorial Algorithms Cornell Notes}
\setCornellUnitType{Chapter}
\setCornellUnitNumber{16}
\setCornellUnitTitle{Permutation Cycles}
\setCornellCompanionLabel{Study and review companion}

\hypersetup{pdftitle={Chapter 16: Permutation Cycles},pdfsubject={Cornell notes},pdfauthor={}}

\begin{document}
\maketitle
\begin{tcolorbox}[colback=Paper,colframe=Ink]{\small\bfseries\color{Teal}CORNELL NOTES}\par{\LARGE\bfseries Chapter 16: The Cycle Structure of a Permutation (CYCLES)}\par\medskip\textbf{Topic:} Cycle counting, sign, tagging, and inversion\hfill\textbf{Date:} \today\par\textbf{Study purpose:} Reuse sign bits to traverse cycles in place.\end{tcolorbox}
\begin{tcolorbox}[colback=Teal!5,colframe=Teal,title=Essential question]How can the cycles of a permutation be counted, tagged, and reversed without an auxiliary visited array?\end{tcolorbox}
\begin{tcolorbox}[colback=white,colframe=Mist,title=Learning objectives]\begin{itemize}\item Decompose a permutation into disjoint cycles.\item Relate cycle count to permutation sign.\item Explain the sign-tagging convention.\item Invert all cycles in place.\item Select CYCLES behavior by option.\end{itemize}\textbf{Prerequisites:} permutation arrays, cycle notation, transpositions, parity.\end{tcolorbox}
\section*{Cornell notes}
\begin{longtable}{@{}Q!{\color{Mist}\vrule width 1pt}A@{}}\cornellhead
\cue{What is the cycle decomposition?}{Following $i,\sigma(i),\sigma^2(i),\ldots$ eventually returns to $i$. Disjoint orbits form cycles and partition $\{1,\ldots,n\}$. If there are $c$ cycles with lengths $n_1,\ldots,n_c$, then $\sum n_j=n$.}
\cue{How is sign obtained?}{A cycle of length $m$ is a product of $m-1$ transpositions. Therefore the whole permutation is a product of $n-c$ transpositions and
\[
\operatorname{sgn}(\sigma)=(-1)^{n-c}.
\]}
\cue{What does tagging mean?}{For each cycle, negate the stored image $\sigma(i)$ at the least index $i$ in that cycle. The negative value marks the point at which a later traversal has completed the cycle and provides a recognizable start without a separate visited array.}
\cue{How does TAG work?}{Scan indices in increasing order. Follow a cycle while temporarily reversing signs of encountered images, reducing the initial cycle count whenever traversal passes to a new member. A final sign-restoration pass leaves negative only the designated tag in each cycle.}
\cue{How is inversion performed?}{After tagging, traverse each cycle and reverse its arrows: if $i_0\mapsto i_1\mapsto\cdots\mapsto i_{m-1}\mapsto i_0$, write the predecessor of every member as its new image. Encountering the negative tag detects the cycle boundary and the signs are removed.}
\cue{What options are supported?}{CYCLES always can calculate cycle count and sign. Option $0$ restores the original permutation; option $1$ leaves the tagged permutation; option $-1$ returns the untagged inverse. The input array is therefore output or workspace depending on the option.}
\cue{Correctness invariant}{Every cycle is entered at its least unprocessed index. Temporary negative signs mark exactly the path already followed. Because cycles are disjoint, processing one cannot corrupt another. Reversing every directed cycle gives the unique inverse mapping.}
\cue{Cautions}{The sign-bit technique assumes values are nonzero labels whose negation is available and unambiguous. Do not apply it directly to unsigned storage. Aliasing or interruption while tags are present exposes a temporarily invalid permutation representation.}
\cue{Retrieval practice}{\begin{enumerate}\item Why is the sign $(-1)^{n-c}$?\item Which entry is tagged per cycle?\item Why is no visited array needed?\item How is a cycle inverted?\item What does option $1$ return?\end{enumerate}}
\end{longtable}
\begin{tcolorbox}[colback=Ink!4,colframe=Ink,title=Cornell summary]
CYCLES exploits the disjoint-cycle structure of a permutation. A cycle of length $m$ contributes $m-1$ transpositions, so a permutation with $c$ cycles has sign $(-1)^{n-c}$. To avoid a separate visited array, the algorithm uses temporary negative signs and leaves one negative tag at the least index of each cycle. Those tags later delimit cycle traversals and make in-place inversion possible by reversing every cycle's arrows. An option selects whether the array is restored, left tagged, or returned as the inverse. The technique is space-efficient but tightly coupled to a signed integer representation.
\end{tcolorbox}
\end{document}
